\subsection{Machine Learning Models}
\label{sec:models}

We implement and evaluate three ML models, chosen to represent different algorithmic families: bagging (Random Forest), boosting (XGBoost), and deep learning (DNN). All implementations use Python with scikit-learn~\cite{ref_sklearn} and TensorFlow/Keras~\cite{ref_tensorflow}.

\subsubsection{Model 1: Random Forest Classifier.}

Random Forest~\cite{ref_rf_breiman} serves as our primary baseline due to its strong performance on tabular data, fast training, and built-in feature importance.

\paragraph{Architecture.}
\begin{itemize}
    \item Number of trees: $B = 100$
    \item Maximum depth: unrestricted (trees grow until leaves are pure or contain fewer than 2 samples)
    \item Split criterion: Gini impurity
    \item Features per split: $\sqrt{d}$ (where $d$ is total features)
    \item Bootstrap sampling: enabled
\end{itemize}

\paragraph{Training procedure.}
\begin{algorithm}[H]
\caption{Random Forest Training for IDS}
\label{alg:rf_train}
\begin{algorithmic}[1]
\Require Training set $\{(\mathbf{x}_i, y_i)\}_{i=1}^{n}$, number of trees $B$
\Ensure Ensemble of $B$ trees $\{h_1, \ldots, h_B\}$
\For{$b = 1$ \textbf{to} $B$}
    \State Draw bootstrap sample $D_b$ of size $n$ from training data
    \State Grow tree $h_b$ on $D_b$:
    \State \quad At each node, select $\sqrt{d}$ random features
    \State \quad Split on the feature/threshold minimizing Gini impurity
    \State \quad Grow until leaf purity or minimum samples reached
\EndFor
\State \textbf{Prediction}: $\hat{y}(\mathbf{x}) = \text{majority\_vote}(h_1(\mathbf{x}), \ldots, h_B(\mathbf{x}))$
\end{algorithmic}
\end{algorithm}

\paragraph{Feature importance.}
Random Forest computes feature importance as the mean decrease in impurity (MDI) across all trees. For feature $j$:
\begin{equation}
    \text{Importance}(j) = \frac{1}{B} \sum_{b=1}^{B} \sum_{t \in T_b} \frac{n_t}{n} \cdot \Delta \text{Gini}(t, j)
\end{equation}
where $T_b$ is the set of nodes in tree $b$ that split on feature $j$, $n_t/n$ is the fraction of samples reaching node $t$, and $\Delta\text{Gini}(t, j)$ is the impurity decrease at that split.

\subsubsection{Model 2: Deep Neural Network (DNN).}

The DNN model leverages deep learning to capture complex, non-linear relationships in the feature space~\cite{ref_dl_ids}.

\paragraph{Architecture.}
\begin{table}[H]
\caption{DNN architecture for intrusion detection.}\label{tab:dnn_arch}
\centering
\begin{tabular}{@{}llll@{}}
\toprule
\textbf{Layer} & \textbf{Type} & \textbf{Units} & \textbf{Details} \\
\midrule
Input & Dense & $d$ & Number of features \\
Hidden 1 & Dense + ReLU & 128 & + BatchNorm + Dropout(0.3) \\
Hidden 2 & Dense + ReLU & 64 & + BatchNorm + Dropout(0.3) \\
Hidden 3 & Dense + ReLU & 32 & + BatchNorm + Dropout(0.2) \\
Output & Dense + Softmax & $K$ & Number of classes \\
\bottomrule
\end{tabular}
\end{table}

\paragraph{Training configuration.}
\begin{itemize}
    \item \textbf{Optimizer}: Adam~\cite{ref_adam} with learning rate $\eta = 0.001$
    \item \textbf{Loss function}: Categorical cross-entropy
    \item \textbf{Batch size}: 256
    \item \textbf{Epochs}: 50 with early stopping (patience = 5, monitoring validation loss)
    \item \textbf{Regularization}: Dropout~\cite{ref_dropout} (0.2--0.3) and Batch Normalization per hidden layer
\end{itemize}

\paragraph{Training procedure.}
\begin{algorithm}[H]
\caption{DNN Training for IDS}
\label{alg:dnn_train}
\begin{algorithmic}[1]
\Require Scaled training data $\{(\mathbf{x}'_i, y_i)\}$, validation split 10\%
\Ensure Trained DNN weights $\theta^*$
\State Initialize weights $\theta$ randomly (He initialization)
\For{epoch $= 1$ \textbf{to} 50}
    \For{each mini-batch $\mathcal{B}$ of size 256}
        \State Forward pass: $\hat{\mathbf{p}} = f_\theta(\mathbf{x}')$ with dropout active
        \State Compute loss: $\mathcal{L} = -\frac{1}{|\mathcal{B}|}\sum_{i \in \mathcal{B}} \sum_{k} y_{ik} \log \hat{p}_{ik}$
        \State Backward pass: compute $\nabla_\theta \mathcal{L}$
        \State Update: $\theta \leftarrow \text{Adam}(\theta, \nabla_\theta \mathcal{L}, \eta=0.001)$
    \EndFor
    \State Evaluate on validation set
    \If{validation loss has not improved for 5 epochs}
        \State \textbf{break} (early stopping)
    \EndIf
\EndFor
\State Restore weights from best validation epoch
\end{algorithmic}
\end{algorithm}

\subsubsection{Model 3: XGBoost Classifier.}

XGBoost~\cite{ref_xgboost} represents the gradient boosting family and is known for consistently winning machine learning competitions on structured data~\cite{ref_ensemble_ids}.

\paragraph{Hyperparameters.}
\begin{itemize}
    \item Number of boosting rounds: 200
    \item Maximum tree depth: 6
    \item Learning rate (shrinkage): $\eta = 0.1$
    \item Column subsampling per tree: 0.8
    \item Row subsampling per tree: 0.8
    \item Regularization: $\lambda = 1.0$ (L2), $\gamma = 0$ (minimum loss reduction)
    \item Objective: \texttt{multi:softprob} (multi-class with probability output)
    \item Early stopping rounds: 10
\end{itemize}

\paragraph{Key advantages for IDS.}
XGBoost offers several properties particularly beneficial for intrusion detection:
\begin{enumerate}
    \item \textbf{Built-in handling of missing values}: Automatically learns the optimal direction for missing splits.
    \item \textbf{Regularization}: L1 and L2 regularization on leaf weights prevents overfitting on imbalanced classes.
    \item \textbf{Feature importance}: Provides gain-based and cover-based importance metrics.
    \item \textbf{Computational efficiency}: Uses histogram-based split finding and cache-aware access patterns for fast training.
\end{enumerate}
